\documentclass[10pt,twocolumn]{article}
\usepackage{hubpaper}
\usepackage{booktabs}
\hubtitle{List Scheduling of Dependent\\ Non-Preemptive Tasks: Priority Rules,\\ Guarantees and an Empirical Comparison}
\begin{document}
\maketitle

\begin{abstract}
\noindent We study the minimisation of the makespan of a set of tasks with precedence constraints on $m$ identical machines without preemption, $P\,|\,\mathrm{prec}\,|\,C_{\max}$. The problem is strongly NP-hard, so we consider list scheduling, for which Graham's bound guarantees a ratio $2-1/m$ whatever the priority rule. We define the classical priorities (height, depth, t-level, b-level, number of successors, duration), a fuzzy tie-breaking rule, and a genetic algorithm that searches over convex combinations of these priorities. On 180 random graphs of 1000 tasks, the b-level rule shortens the makespan by $2.43\%$ on average relative to the height rule (standard error $0.30$); the fuzzy tie-break and the genetic algorithm give no significant further gain.
\end{abstract}

\section{Problem}
An instance is a directed acyclic graph $G=(V,E)$, $|V|=n$, with durations $p_j>0$, and $m$ identical machines. A schedule assigns start times $S_j\ge0$ and machines such that (i) $S_j\ge S_i+p_i$ for all $(i,j)\in E$; (ii) a task runs uninterrupted on one machine during $[S_j,S_j+p_j)$; (iii) at most $m$ tasks run at any time. The objective is $C_{\max}=\max_j(S_j+p_j)$. In the notation of Graham et al.~\cite{graham1979} this is $P\,|\,\mathrm{prec}\,|\,C_{\max}$; it is NP-hard already for unit durations~\cite{ullman1975} and strongly NP-hard in general~\cite{lenstra1978}, which rules out exact methods on graphs with thousands of tasks; see~\cite{kwok1999} for a survey of heuristics.

\textbf{Lower bound.} Let the length of a path be the sum of its durations and $\mathrm{CP}$ the length of a longest (critical) path. Any schedule satisfies
\begin{equation}\label{eq:lb}
C^*_{\max}\ge\mathrm{LB}=\max\Big(\mathrm{CP},\;\tfrac1m\textstyle\sum_jp_j\Big),
\end{equation}
since the tasks of a path run sequentially and the total work is shared by $m$ machines.

\section{List scheduling and Graham's bound}
Given a priority $\pi:V\to\R$, list scheduling is the event-driven rule: whenever a machine is idle and the set of \emph{ready} tasks (all predecessors finished) is non-empty, start the ready task of highest priority.
\begin{theorem}[Graham~\cite{graham1966,graham1969}]
For every priority $\pi$, the list schedule satisfies
$C_{\max}\le\frac1m\sum_jp_j+\big(1-\frac1m\big)\mathrm{CP}\le\big(2-\frac1m\big)C^*_{\max}.$
\end{theorem}
\begin{proof}
Let $j_k$ be a task finishing last and $C_j=S_j+p_j$. Given $j_\ell$, let $j_{\ell-1}$ be a predecessor of $j_\ell$ finishing last. On $[C_{j_{\ell-1}},S_{j_\ell})$, $j_\ell$ is ready, so no machine is idle there (the rule never idles a machine while a task is ready). Stop at a task $j_1$ without predecessors, which is ready on $[0,S_{j_1})$. Hence at every instant of $[0,C_{\max})$ not covered by a chain task, all $m$ machines are busy. Counting busy time, $mC_{\max}\le\sum_jp_j+(m-1)\sum_\ell p_{j_\ell}\le\sum_jp_j+(m-1)\mathrm{CP}$. Conclude with \eqref{eq:lb}.
\end{proof}
The bound is tight in the worst case and list scheduling exhibits anomalies: adding a machine or shortening a task may increase $C_{\max}$~\cite{graham1969}. The choice of $\pi$ therefore matters only on average, which is what we measure.

\section{Priority rules}
With $\mathrm{pred}(j)$, $\mathrm{succ}(j)$ the immediate neighbours, the rules are computed by one pass in topological (resp.\ reverse topological) order, in $O(n+|E|)$:
\begin{align*}
\mathrm{height}(j)&=\textstyle\max_{i\in\mathrm{pred}(j)}\mathrm{height}(i)+1,\\
\mathrm{depth}(j)&=\textstyle\max_{k\in\mathrm{succ}(j)}\mathrm{depth}(k)+1,\\
t\ell(j)&=p_j+\textstyle\max_{i\in\mathrm{pred}(j)}t\ell(i),\\
b\ell(j)&=p_j+\textstyle\max_{k\in\mathrm{succ}(j)}b\ell(k),
\end{align*}
with $\max\emptyset=0$; to these we add $|\mathrm{succ}(j)|$ and $p_j$. The b-level is the length of the longest path starting at $j$: even with infinitely many machines, nothing can finish before $S_j+b\ell(j)$, so $C_{\max}\ge\max_j(S_j+b\ell(j))$ and the task with the largest b-level is the one whose delay is most likely to delay the whole schedule; $\max_jb\ell(j)=\mathrm{CP}$. Depth is the b-level with unit durations. For in-trees with unit durations, the highest-level-first rule is optimal~\cite{hu1961}, and Coffman--Graham is optimal for $m=2$~\cite{coffman1972}; b-level is the natural extension to arbitrary durations and is the core of HEFT~\cite{topcuoglu2002}.

\textbf{Fuzzy tie-breaking.} Among tasks of equal b-level, a Mamdani-type rule base~\cite{zadeh1965} on height and duration, with triangular memberships $\mu(x;a,b,c)=\max\!\big(0,\min(\frac{x-a}{b-a},\frac{c-x}{c-b})\big)$, the $\min$ t-norm and centroid defuzzification, gives a secondary key. It can only permute tasks inside classes of equal b-level, which are small when durations are continuous; its effect is therefore expected to be small.

\section{Genetic search over priorities}
Let $f_1,\dots,f_6$ be the six features above, normalised to $[0,1]$. A chromosome encodes $w\in\{0,\dots,255\}^6$ (48 bits) and the priority $\pi_w(j)=\sum_kw_kf_k(j)$; its fitness is the makespan of the corresponding list schedule. Selection keeps the best chromosomes, single-point crossover (rate $0.5$) and bit-flip mutation (rate $0.2$) produce offspring~\cite{holland1975}. Every pure rule is a point of the search space ($w=e_k$), so if the population is seeded with the six pure rules and the best chromosome is always kept (elitism), the best makespan found is non-increasing and never worse than the best pure rule; without this seeding, nothing guarantees it.

\section{Experiments}
Graphs are random DAGs, $m$ is drawn between $2\%$ and $5\%$ of $n$, and performance is $R_A=100\,(1-C^A_{\max}/C^{\mathrm{height}}_{\max})$ (positive is better). Table~\ref{tab:res} reports mean and standard deviation over 180 graphs, with standard error $s/\sqrt{180}$.
\begin{table}[h]\centering\small
\begin{tabular}{lrrr}\toprule
& \multicolumn{2}{c}{$n=1000$} & $n=50$\\ \cmidrule(lr){2-3}
Rule & mean (sd) & s.e. & mean (sd)\\\midrule
b-level & $2.43$ ($4.01$) & $0.30$ & $3.11$ ($3.39$)\\
b-level + fuzzy & $2.43$ ($4.00$) & $0.30$ & $3.13$ ($3.40$)\\
depth & $1.79$ ($3.05$) & $0.23$ & $1.11$ ($2.77$)\\
successors & $-0.14$ ($1.75$) & $0.13$ & $0.27$ ($2.41$)\\
duration & $-0.29$ ($1.35$) & $0.10$ & $0.61$ ($3.34$)\\
t-level & $-1.08$ ($2.18$) & $0.16$ & $-2.34$ ($4.00$)\\
genetic & -- & -- & $0.25$ ($2.28$)\\\bottomrule
\end{tabular}
\caption{Improvement over the height rule (\%).}\label{tab:res}
\end{table}

\textbf{Findings.} (1) b-level dominates: its mean gain is about $8$ standard errors above zero, and depth, its unit-duration version, is second, which confirms that the information carried by the remaining critical path is what matters. (2) t-level is worse than the baseline: it favours tasks far from the sources, i.e.\ it looks backwards when the makespan is determined by what remains. (3) The fuzzy tie-break does not differ from b-level within noise, as predicted. (4) The genetic algorithm, with population 30 and 20 generations ($600$ fitness evaluations), did not beat b-level on $n=50$ and was by far the slowest; on such small graphs most priorities give the same schedule, so selection has no signal. All rules except the genetic algorithm run within $0.1$\,s of each other.

\textbf{Conclusion.} b-level list scheduling is, at equal cost, the best of the tested rules, with Graham's $2-1/m$ guarantee. A seeded, elitist genetic search, or a hybrid with local search, is the natural next step to obtain further improvement on large graphs.

\footnotesize
\begin{thebibliography}{99}\setlength{\itemsep}{0pt}
\bibitem{coffman1972} E.~Coffman, R.~Graham. Optimal scheduling for two-processor systems. \emph{Acta Informatica} 1:200--213, 1972.
\bibitem{graham1966} R.~Graham. Bounds for certain multiprocessing anomalies. \emph{Bell Syst.\ Tech.\ J.} 45(9):1563--1581, 1966.
\bibitem{graham1969} R.~Graham. Bounds on multiprocessing timing anomalies. \emph{SIAM J.\ Appl.\ Math.} 17(2):416--429, 1969.
\bibitem{graham1979} R.~Graham, E.~Lawler, J.~Lenstra, A.~Rinnooy~Kan. Optimization and approximation in deterministic sequencing and scheduling: a survey. \emph{Ann.\ Discrete Math.} 5:287--326, 1979.
\bibitem{holland1975} J.~Holland. \emph{Adaptation in Natural and Artificial Systems}. Univ.\ of Michigan Press, 1975.
\bibitem{hu1961} T.~C.~Hu. Parallel sequencing and assembly line problems. \emph{Operations Research} 9(6):841--848, 1961.
\bibitem{kwok1999} Y.-K.~Kwok, I.~Ahmad. Static scheduling algorithms for allocating directed task graphs to multiprocessors. \emph{ACM Comput.\ Surv.} 31(4):406--471, 1999.
\bibitem{lenstra1978} J.~Lenstra, A.~Rinnooy~Kan. Complexity of scheduling under precedence constraints. \emph{Operations Research} 26(1):22--35, 1978.
\bibitem{topcuoglu2002} H.~Topcuoglu, S.~Hariri, M.-Y.~Wu. Performance-effective and low-complexity task scheduling for heterogeneous computing. \emph{IEEE TPDS} 13(3):260--274, 2002.
\bibitem{ullman1975} J.~Ullman. NP-complete scheduling problems. \emph{J.\ Comput.\ Syst.\ Sci.} 10(3):384--393, 1975.
\bibitem{zadeh1965} L.~Zadeh. Fuzzy sets. \emph{Information and Control} 8(3):338--353, 1965.
\end{thebibliography}
\end{document}
